\begin{trapbox}

	\begin{description}[style=nextline, leftmargin=2.8em, labelsep=0.5em, itemsep=0.35em]
		\item[\textbf{\textcolor{CTrap}{易错点一（忽略定义域对称）}}]
			错误做法：只要 $f(a+x)=f(a-x)$ 成立就认为图象对称，不检查定义域。
		\item[\textbf{\textcolor{CTrap}{正确理解（定义域前提）：}}]
			定义域必须关于 $x=a$ 对称，否则即使等式成立也不能保证整体对称。
		\item[\textbf{\textcolor{CTrap}{错因分析（对称点缺失）：}}]
			对称定义要求每个点都有对称点，定义域不对称则部分点无对称点。
	\end{description}

	\medskip
	\begin{description}[style=nextline, leftmargin=2.8em, labelsep=0.5em, itemsep=0.35em]
		\item[\textbf{\textcolor{CTrap}{易错点二（混淆对称公式）}}]
			错误做法：将 $f(a+x)=f(a-x)$ 混淆为 $f(x)=f(a-x)$ 或 $f(x)=f(x-a)$。
		\item[\textbf{\textcolor{CTrap}{正确理解（正确等价形式）：}}]
			正确等价形式为 $f(a+x)=f(a-x)$ 或 $f(x)=f(2a-x)$。
		\item[\textbf{\textcolor{CTrap}{错因分析（变量替换出错）：}}]
			变量替换混淆，导致对称轴判断错误。
	\end{description}

	\medskip
	\begin{description}[style=nextline, leftmargin=2.8em, labelsep=0.5em, itemsep=0.35em]
		\item[\textbf{\textcolor{CTrap}{易错点三（片面验证）}}]
			错误做法：取一个或几个 $x$ 值验证等式成立即断言对称。
		\item[\textbf{\textcolor{CTrap}{正确理解（整体性要求）：}}]
			必须对定义域内所有允许的 $x$ 恒成立。
		\item[\textbf{\textcolor{CTrap}{错因分析（整体误判）：}}]
			对称性是整体性质，不能由个别点断定。
	\end{description}

\end{trapbox}
